\scanpage{53}
% @ lines 1–2 / text
De plus, cet argument montre aussi que $J$ est uniquement déterminé comme $\operatorname{Ker}(A\to\operatorname{End}(\underline L))$.
% @ lines 3–8 / text
De plus, sous les conditions équivalentes a) b) c) \deleted{on aura évidemment $\det|\underline J=0$ donc $J$ \unsure{se}.} soit loc. $u$ une base de $\underline J$, \deleted{\unsure{donc est soit|donc est sur} de la forme $\underline L\otimes$} de sorte que $L=\operatorname{Ker}u$, \unsure{soit|donc} $M=\operatorname{Im}u$, c'est un ss-fibré vectoriel inv. tel que \deleted{$L^\perp\otimes M$} $L^\vee\otimes M=\mathcal O_Su=\underline J\subset A$, d'où d).
% @ lines 8–13 / text
D'ailleurs, si $M$ est un ss-fibré inv. de $V$ tel que $L^\vee\otimes M\subset A$, évidemt $L^\vee\otimes M$ est un ss-fibré inv. de $A$ nul sur $L$, donc c'est $J$, ce qui montre en $=$ temps que $M$ est uniquement déterminé comme $\operatorname{Im}u$ ($u$ base loc$^{\mathrm{le}}$ de $J$).
% @ lines 14–17 / text
Reste à prouver d)$\longrightarrow$b), or \unsure{d)} implique $A=\mathcal O_S+J$ (où on pose $J=L^\perp\otimes M$) \deleted{et} $J$ étant nul sur \deleted{\unsure{$\mathcal O$}} $\underline L$ à fortiori l'invarie cqfd.

% @ line 18, overwritten heading / deletion
\deleted{\uline{Corollaire}}

% @ lines 19–23 / text
Se donner un $A$ (ss fibré vectoriel de rg $2$ de $M$ contenant $1$\inserted{i.e.} $\mathcal O_S.1$ -- un tel $A$ est néc une ss-alg comm !) ou se donner $A^\perp=H\subset\widetilde{\unsure{\mathfrak g|g}}$, un sous fibré vectoriel de rang $2$ \inserted{i.e. corg $1$,} de $\widetilde{\mathfrak g}$, revient au $=$
% @ line 24, lower edge / display
\[
 \operatorname{AlgSpec}(M)\simeq
 \underbrace{\mathbb P(\widetilde{\mathfrak g})\;(\simeq\check{\mathbb P}(\mathfrak g))}_{\check P}
 \simeq\left\{\begin{array}{l}
 \inserted{\text{schéma des}}\\[-2pt]
 \text{fibrés en droites)}\\
 \text{dans }P\mathbin{:}\check{\mathbb P}(\widetilde{\mathfrak g})=\mathbb P(\mathfrak g)
 \end{array}\right.
 \tag*{(21)}
\]
% archive sheet number: 52
